\section{Equivalence and equisatisfiability}
\label{sec:automated-reasoning-02}

\begin{definition}[Equisatisfiability]
Formulas \(\varphi\) and \(\psi\), possibly in different signatures, are
\emph{equisatisfiable} iff the two formulas are either both satisfiable
or both unsatisfiable.
\end{definition}

Logical equivalence is stronger: it compares truth under corresponding
valuations, whereas equisatisfiability compares only existence of a satisfying
valuation.  Tseitin clausification \cite{Tseitin1968} introduces a fresh
variable for each
selected subformula and constrains that variable to have the same truth value
as the subformula.  If \(p_\varphi\) names the complete input formula
\(\varphi\), the output also contains the unit clause \(\set{p_\varphi}\),
which requires the root variable to be true.  Call the resulting finite clause
set \(S_\varphi\).  The conjunction of the
\(\operatorname{form}_{\CPL}\)-images in any enumeration of
\(S_\varphi\) is equisatisfiable with \(\varphi\) in the expanded signature,
and
\(S_\varphi\) can have size linear in the size of \(\varphi\).  The output is
not literally the same formula and need not be logically equivalent in the
original signature.

\begin{example}[Tseitin constraint]
For a subformula \(\theta=(P\land Q)\), introduce fresh \(p_\theta\) and add
the finite set of clauses
\[
  S_\theta=
  \set{
    \set{\neg p_\theta,P},
    \set{\neg p_\theta,Q},
    \set{p_\theta,\neg P,\neg Q}
  }.
\]
Conjoining the three \(\operatorname{form}_{\CPL}\)-images in the displayed
clause order gives a formula logically equivalent to
\[
  (\neg p_\theta\lor P)\land
  (\neg p_\theta\lor Q)\land
  (p_\theta\lor\neg P\lor\neg Q),
\]
which expresses \(p_\theta\leftrightarrow(P\land Q)\).
Every valuation of \(P,Q\) extends uniquely to \(p_\theta\) while satisfying
\(S_\theta\).  If \(\theta\) is the complete input, the Tseitin output is
\(S_\theta\cup\set{\set{p_\theta}}\); the added unit clause requires the
extension to make \(\theta\) true.
\end{example}

